\documentclass[11pt]{article}

\usepackage[margin=25mm]{geometry}

\usepackage{amsmath,amssymb,fontspec,unicode-math}

\IfFontExistsTF{DejaVu Serif}{\setmainfont{DejaVu Serif}}{\setmainfont{Latin Modern Roman}}

\IfFontExistsTF{Latin Modern Math}{\setmathfont{Latin Modern Math}}{\setmathfont{DejaVu Math TeX Gyre}}

\usepackage{xurl,hyperref}

\setlength{\parindent}{0pt}\setlength{\parskip}{6pt}

\begin{document}

\title{Don’t Just Chase “Highlighted Tokens” in MLLMs: Revisiting Visual Holistic Context Retention}

\author{ELVIA preliminary AI analysis}\date{21 October 2026}\maketitle

Preliminary assessment: Minor.\par

Independent scientific review remains pending. Author adoption has not been requested.

The proportional allocator does not inherit a greedy approximation guarantee. Additional semantic-pruning proof and public selection/quota paths need correction, without a demonstrated failure of the reported visual-language results.

\section*{Review scope}

Latest-version, claim-focused review of the cited author-paper passages and applicable current author-release paths. The stated counterexamples and scope controls were checked analytically; no author program, old verifier, training run or benchmark was executed.

\section*{Problem 1: Proportional allocation is not the guaranteed greedy algorithm}

Source: Section 4.1, PDF p.6, Eq.(4); Appendix C, Corollary 1, PDF p.31; proof C.3, PDF p.32

Assessment: supported; impact: minor

\subsection*{Source claim (paraphrase)}

Appendix C attaches a \(1-1/e\) guarantee to the allocation discussion, whereas Eq.(4) uses proportional crop scores.

\subsection*{Analysis}

For two crops with \(100\) constant-score tokens each, scores \(1\) and \(1/99\) yield quotas \((99,1)\) and log utility \(0\). Feasible quotas \((50,50)\) give \(\log(2500/99)>0\). Thus the claimed positive multiplicative guarantee cannot be attached to that proportional allocation. A separately specified marginal-gain greedy procedure is not refuted.

\subsection*{Finding and scope}

The supporting algorithm-to-guarantee attribution is invalid; empirical usefulness of proportional quotas remains open.

\subsection*{Proposed correction}

For positive sorted scores, initialize one token per crop and allocate remaining feasible tokens by maximal log-prefix marginal gain. Its diminishing marginal sequences justify the appropriate separable allocation argument.

An explicit local correction is described; author adoption and a corresponding published erratum are not established.

\section*{Problem 2: Token coverage does not by itself control pruning normalization}

Source: Appendix C, PDF p.31, Assumption 1, Lemma C.1, Theorem C.1 and proof C.2

Assessment: supported; impact: minor

\subsection*{Source claim (paraphrase)}

Appendix C passes from token similarity/coverage assumptions to a transformer-output preservation bound.

\subsection*{Analysis}

With zero queries, attention on \([e_1,e_1,e_2]\) gives \((2/3,1/3)\), while pruning one duplicate gives \((1/2,1/2)\). Perfect coverage of the removed duplicate does not preserve multiplicity or normalization. A fixed-shape Lipschitz argument needs a common lifted representation and a separate deletion bound. The unspecified remainder may absorb this example.

\subsection*{Finding and scope}

A proof bridge is missing; the full numerical theorem is not refuted by this control.

\section*{Problem 3: Selection uses raw attention instead of the computed holistic score}

Source: Section 4.1, PDF p.6, Eqs.(3)–(5); llava/model/utils.py, lines 22–164

Assessment: supported; impact: minor

\subsection*{Source claim (paraphrase)}

Equation (5) selects holistic scores; `llava/model/utils.py`, `HoloV`, computes those scores but takes top entries from `patch\_attn`. The mean-normalized crop term also largely cancels.

\subsection*{Analysis}

For duplicated \(e_1\) and one \(e_2\), attentions \((0.333,0.333,0.334)\) can rank the \(e_2\) token highest by attention while holistic scores favor \(e_1\). Away from degeneracy, crop normalization makes the diversity contribution nearly a common factor, with rounding boundaries still relevant. Repository issues \#13 and \#15 already discuss parameter insensitivity.

\subsection*{Finding and scope}

The released path differs from the score definition; historical parameter effects are not newly quantified.

\subsection*{Proposed correction}

Use the intended holistic score for selection and an explicitly defined crop aggregation that does not unintentionally normalize away the desired variation. No recovered benchmark result is established.

An explicit local correction is described; author adoption and a corresponding published erratum are not established.

\section*{Problem 4: Overflow quotas produce padding instead of feasible real tokens}

Source: Section 4.1, PDF p.6, quota redistribution; llava/model/utils.py, lines 100–157

Assessment: supported; impact: minor

\subsection*{Source claim (paraphrase)}

Section 4.1 specifies capacity caps and redistribution; the released function floors quotas and later pads short outputs without redistributing overflow.

\subsection*{Analysis}

Two crops of capacity \(2\), weights \((0.98,0.02)\) and budget \(3\) produce nominal quotas \((3,0)\). The path selects only two real tokens and pads a zero, although \((2,1)\) is feasible.

\subsection*{Finding and scope}

The overflow branch violates the stated quota rule. Nonoverflow behavior is unaffected.

\subsection*{Proposed correction}

Cap initial quotas by capacity, then repeatedly allocate remaining tokens to an eligible crop under a declared priority rule until the feasible budget is met.

An explicit local correction is described; author adoption and a corresponding published erratum are not established.

\section*{Limits}

Preliminary AI assessment; independent human scientific adjudication remains pending.\par

An analytic counterexample or a release discrepancy does not by itself remeasure the paper’s reported experiments.\par

\section*{Source identity}

NeurIPS 2025 official proceedings PDF; compared with current arXiv:2510.02912v2

\url{https://proceedings.neurips.cc/paper_files/paper/2025/file/38b445e98823a02479582302ee3aa8e4-Paper-Conference.pdf}

Reviewed PDF SHA-256: \nolinkurl{158d2655adeb5bdf25623b5b8a1166339cad1dcc76d243fe29e0d52eb9439672}

The selected author/publisher source is the same work as the original case, matched by exact title, author identity and official version channel. The exact selected PDF SHA-256 was recomputed in the latest-version pass for Issue \#27. The PDF is a retained author/publisher copy; its original retrieval time is stated only when supported by matching acquisition metadata. File modification time was not substituted.

\end{document}
